% BHU PYQ -- Manually transcribed & error-corrected
\documentclass[11pt,a4paper]{article}

\usepackage[a4paper,top=2.5cm,bottom=2.5cm,left=2.8cm,right=2.8cm,headheight=14pt]{geometry}
\usepackage{amsmath,amssymb}
\usepackage[T1]{fontenc}
\usepackage[utf8]{inputenc}
\usepackage{lmodern}
\usepackage{microtype}
\usepackage{fancyhdr}
\usepackage{enumitem}
\usepackage{hyperref}
\usepackage{lastpage}
\usepackage{parskip}

\hypersetup{colorlinks=true,urlcolor=black,linkcolor=black,
  pdftitle={BPE-601: Topics in Modern Physics -- BHU PYQ},pdfauthor={Banaras Hindu University}}

\pagestyle{fancy}\fancyhf{}
\renewcommand{\headrulewidth}{0.4pt}\renewcommand{\footrulewidth}{0.4pt}
\fancyhead[L]{\small   \textit{Banaras Hindu University}}
\fancyhead[R]{\small   \textit{BPE-601: Topics in Modern Physics}}
\fancyfoot[C]{\small Page  \thepage\ of \pageref{LastPage}}


\newlist{parts}{enumerate}{2}
\setlist[parts,1]{label=(\alph*),leftmargin=*,itemsep=5pt,topsep=3pt}
\setlist[parts,2]{label=(\roman*),leftmargin=*,itemsep=3pt,topsep=2pt}

\newcommand{\pts}[1]{\hfill   \textit{[#1]}}


\begin{document}


\begin{center}
  {\large\bfseries BANARAS HINDU UNIVERSITY}\\[2pt]
  {
\normalsize B.Sc. (Hons.) Semester VI Examination, 2013-14}\\[6pt]
  \rule{0.8\linewidth}{0.5pt}\\[6pt]
  {\large\bfseries Paper No. BPE-601: Topics in Modern Physics}\\[4pt]
  {
\normalsize Time: 3 Hours\hfill Maximum Marks: 70}
\end{center}
\rule{\linewidth}{0.4pt}
\smallskip



\noindent   \textit{Instructions: Answer all questions. Symbols have their usual meanings.}
\bigskip




\noindent   \textbf{Question 1.} Answer any   \textbf{five} of the following: \pts{5 $ \times$ 4 = 20}

\begin{parts}
  \item State and explain the Cosmological Principle.
  \item What are degenerate stars? Explain briefly with an example.
  \item Draw neatly the Hertzsprung-Russell (H-R) diagram and explain it.
  \item Write down the Lorentz transformations for energy and momentum when the boost is along the $x$-axis.
  \item What are Stokes and anti-Stokes lines in Raman scattering?
  \item What are gauge transformations? Show that the electric and magnetic fields remain invariant under gauge transformations.
\end{parts}

\medskip\rule{\linewidth}{0.2pt}




\noindent   \textbf{Question 2.}

\begin{parts}
  \item Derive an expression for the kinetic energy of the recoiled electron in Compton scattering. \pts{8}
  \item What is the physical meaning of Compton wavelength? Justify your answer. \pts{4.5}
\end{parts}
\hfill   \textbf{OR}

\begin{parts}
  \item Derive the Hamiltonian for a charged particle moving under the influence of an electromagnetic field. \pts{8}
  \item Explain briefly the working of a p-n junction diode as a radiation detector. \pts{4.5}
\end{parts}

\medskip\rule{\linewidth}{0.2pt}




\noindent   \textbf{Question 3.}

\begin{parts}
  \item Discuss the Big Bang nucleosynthesis. \pts{6}
  \item Explain the differences between Big Bang nucleosynthesis and stellar nucleosynthesis. \pts{6.5}
\end{parts}
\hfill   \textbf{OR}

\begin{parts}
  \item What are source-free Maxwell's equations? Write these equations in four-vector notation. \pts{8}
  \item Derive the relationship $m = \frac{m_0}{\sqrt{1 - v^2/c^2}}$, where $m_0$ is the rest mass of the particle, $m$ is the relativistic mass, and $v$ is the velocity of the particle. \pts{4.5}
\end{parts}

\medskip\rule{\linewidth}{0.2pt}




\noindent   \textbf{Question 4.}

\begin{parts}
  \item Derive a condition for the hydrostatic equilibrium of a mass of gas contracting under gravity. \pts{4.5}
  \item Explain the process of generation of energy by the Sun. Discuss its evolution and ultimate fate (i.e., its death). \pts{8}
\end{parts}
\hfill   \textbf{OR}

\begin{parts}
  \item State and explain the significance of Hubble's law. \pts{4}
  \item What is cosmic microwave background? Explain its origin. \pts{4}
  \item What are the various structures in the visible universe? What does it mean by large-scale smoothness of the universe? \pts{4.5}
\end{parts}

\vfill
\rule{\linewidth}{0.4pt}

\begin{center}\small
  \href{https://bhu-exam.vercel.app/}{Click here to visit our website}\\[4pt]
  \href{https://me-aryan.vercel.app/}{Click here to visit our contributor}\\[4pt]
  \href{https://quashan-me.vercel.app/}{Click here to visit our contributor}
\end{center}

\end{document}
